% Generated draft: requires mathematical and source review.
\documentclass[11pt]{article}
\usepackage[margin=25mm]{geometry}
\usepackage{fontspec}
\setmainfont{Latin Modern Roman}
\usepackage{amsmath,amssymb,mathtools,unicode-math}
\setmathfont{Latin Modern Math}
\usepackage{enumitem,needspace,xurl,hyperref}
\setlength{\parindent}{0pt}
\setlength{\parskip}{5pt}
\setlength{\emergencystretch}{3em}
\newcommand{\N}{\mathbb{N}}
\newcommand{\Z}{\mathbb{Z}}
\newcommand{\Q}{\mathbb{Q}}
\newcommand{\R}{\mathbb{R}}
\newcommand{\C}{\mathbb{C}}
\newcommand{\F}{\mathbb{F}}
\newcommand{\A}{\mathbb{A}}
\newcommand{\PP}{\mathbb{P}}
\newcommand{\E}{\mathbb{E}}
\newcommand{\eps}{\varepsilon}
\newcommand{\dd}{\,\mathrm{d}}
\begin{document}
\section*{UnsolvedMath ID: 30005006}
\textbf{Source problem code:} catalogue entry
\subsection{Definitions and mathematical statement}
\textbf{Unrecoverable statement.} The supplied record for permanent UnsolvedMath ID \(30005006\), source code \texttt{catalogue entry}, does not contain a recoverable mathematical assertion. Its statement begins with the isolated word ``conjecture'', followed by a contents-page reference to page 393 and entries for several workshop contributions. The supplied title, ``Unspecified Polynomial Interpolation Conjecture'', does not determine which conjecture is intended.

No mathematical objects, definitions, hypotheses, domains, boundary conditions, quantifiers, or conclusion can be recovered from this fragment. Accordingly, no precise conjecture can be stated from the supplied evidence. The later contents entry entitled ``A Conjecture on Bernstein Basis Polynomials'' does not establish that it is the intended target.

The record has \texttt{statement\_status: unrecoverable} and no clean statement. Its supplied verification note says that the official report identifies the extract as a contents fragment without a recoverable standalone problem; this is a supplied source-verification assessment, not an independent verification here. The supplied \texttt{open} label reflects the curation rule for an unresolved extraction and does not establish that any identified mathematical conjecture remains open.
\subsection{Short English statement}
The supplied text is a truncated workshop table of contents, so it does not identify a mathematical problem precisely enough to formalize. The intended conjecture and its assumptions must be recovered before a mathematical statement can be drafted. The supplied open status is a curation label.
\subsection{Sources}
\begin{itemize}
\item Local UnsolvedMath registry, supplied record with permanent ID \(30005006\), exact source code \texttt{catalogue entry}, title ``Unspecified Polynomial Interpolation Conjecture'', in the collection ``Oberwolfach Reports Open Problems''. No proposer is supplied; the proposed-year field is 2022.
\item Supplied report citation: \emph{Mini-Workshop: Interpolation, Approximation, and Algebra} (2023), \emph{Oberwolfach Reports}, pp. 383--428, DOI \texttt{10.4171/owr/2022/7}. Supplied source URL and literature-source link, labelled ``OWR interpolation report'': \url{https://doi.org/10.4171/owr/2022/7}. The citation's 2023 date and the record's 2022 proposed year are retained without reconciliation.
\item Attribution within the supplied contents fragment: Tomas Sauer (joint with Jesus Carnicer and Carmen Godés), ``Berzolari, Radon, H-bases, syzygies and maximal lines'', p. 396; Evelyne Hubert and Michael Singer, ``Sparse Interpolation in Terms of Multivariate Chebyshev Polynomials'', p. 397; Thomas Fink, ``Analysis of the geometry of edges with the Taylorlet transform'', p. 401; Jesús Carnicer (joint with Carmen Godés), ``GC sets with few maximal lines'', p. 403; Hakop Hakopian, ``On n-independent, n-correct, and GCn sets'', p. 405; Alexandra Seceleanu, ``Higher order interpolation through the singular locus of a reflection arrangement'', p. 408; Boris Shekhtman, ``On Polynomial Interpolation on Arbitrary Varieties'', p. 411; Larry L. Schumaker, ``A Conjecture on Bernstein Basis Polynomials'', p. 413; Kristian Ranestad, ``Adjoints, a short introduction with a view to interpolation and Calabi-Yau varieties'', p. 413; Yuan Xu, ``Polynomial interpolation and cubature rules'', p. 415; Nelly Villamizar (joint with Michael DiPasquale), ``On the dimension of trivariate spline spaces'', p. 417; Michael Floater, ``GBCs Part 2'', p. 419; Tatyana Sorokina, ``Harmonic splines'', p. 420. These are contents-entry attributions, not established authorship of the unidentified conjecture ending on page 393.
\item Supplied literature assessment dated exactly 2026-08-22: the extracted record is a truncated contents fragment ending in ``conjecture'' and does not identify a mathematical claim; under the stated curation rule, the unresolved record is classified as open. The listed source checked is the OWR interpolation report linked above. This is supplied dated curation, not independently verified literature research. The background explicitly requires independent verification of time-sensitive status and cited claims.
\end{itemize}
\end{document}
